\begin{tikzpicture}[>=Latex]

    % --- Definición de colores ---
    \colorlet{wallcolor}{gray!50}
    \colorlet{fluidcolor}{cyan!60!blue}
    \colorlet{pistoncolor}{orange!50!brown}
    \colorlet{f1color}{cyan!70!blue}
    \colorlet{f2color}{blue!40!darkgray}
    \colorlet{carred}{red!80!black}

    % --- 1. Estructura del recipiente (Paredes) ---
    % Pared izquierda
    \fill[wallcolor] (0, 0) rectangle (1, 5);
    % Pared inferior
    \fill[wallcolor] (1, 0) rectangle (9, 1);
    % Pared derecha
    \fill[wallcolor] (9, 0) rectangle (10, 5);
    % Pared central (isla interior)
    \fill[wallcolor] (2, 2) rectangle (6, 5);

    % Bordes de las paredes para dar volumen/definición
    \draw[thick, gray!80!black] (0,5) -- (0,0) -- (10,0) -- (10,5) -- (9,5) -- (9,1) -- (1,1) -- (1,5) -- cycle;
    \draw[thick, gray!80!black] (2,5) -- (2,2) -- (6,2) -- (6,5) -- cycle;

    % --- 2. Líquido (Azul) ---
    % Sigue la forma interior de las paredes
    \fill[fluidcolor] (1, 1) -- (1, 4.5) -- (2, 4.5) -- (2, 2) -- (6, 2) -- (6, 4.0) -- (9, 4.0) -- (9, 1) -- cycle;

    % --- 3. Pistones ---
    % Pistón pequeño (A1)
    \fill[pistoncolor] (1, 4.5) rectangle (2, 4.8);
    \draw[thick, orange!30!black] (1, 4.5) rectangle (2, 4.8);

    % Pistón grande (A2)
    % Base cilíndrica interior
    \fill[pistoncolor] (6.2, 4.0) rectangle (8.8, 5.2);
    \draw[thick, orange!30!black] (6.2, 4.0) rectangle (8.8, 5.2);
    % Plataforma superior exterior
    \fill[pistoncolor!90!black] (5.6, 5.2) rectangle (9.4, 5.5);
    \draw[thick, orange!30!black] (5.6, 5.2) rectangle (9.4, 5.5);

    % --- 4. Vehículo (Coche rojo simplificado) ---
    \begin{scope}[shift={(7.5, 5.5)}]
        % Carrocería principal
        \fill[carred, rounded corners=2pt] (-1.3, 0.2) rectangle (1.3, 0.8);
        % Cabina
        \fill[carred, rounded corners=2pt] (-0.8, 0.8) -- (-0.5, 1.4) -- (0.5, 1.4) -- (1.1, 0.8) -- cycle;
        
        % Ruedas
        \fill[black] (-0.8, 0.2) circle (0.25);
        \fill[black] (0.8, 0.2) circle (0.25);
        \fill[gray!40] (-0.8, 0.2) circle (0.1); % Tapón izquierdo
        \fill[gray!40] (0.8, 0.2) circle (0.1);  % Tapón derecho
        
        % Ventanas
        \fill[cyan!20, rounded corners=1pt] (-0.6, 0.85) -- (-0.4, 1.3) -- (-0.05, 1.3) -- (-0.05, 0.85) -- cycle;
        \fill[cyan!20, rounded corners=1pt] (0.05, 0.85) -- (0.05, 1.3) -- (0.4, 1.3) -- (0.8, 0.85) -- cycle;
        
        % Faros
        \fill[yellow!80!white] (1.2, 0.5) circle (0.08); % Delantero
        \fill[red!50!black] (-1.2, 0.45) rectangle (-1.3, 0.65); % Trasero
        
        % Antena
        \draw[thin] (-0.3, 1.4) -- (-0.6, 1.8);
    \end{scope}

    % --- 5. Flechas y Etiquetas ---

    % Fuerza F1 (entrante, izquierda)
    \draw[->, line width=2.5pt, f1color] (1.5, 6.8) -- (1.5, 4.9);
    \node[] at (1.5, 7.2) {$\vec{\mathbf{F}}_1$};

    % Fuerza F2 (saliente, derecha)
    \draw[->, line width=3pt, f2color] (7.5, -0.5) -- (7.5, 2.5);
    \node[text=white] at (7.9, 1.8) {$\vec{\mathbf{F}}_2$};

    % Etiqueta A1
    \node[] (A1) at (3.0, 4.65) {$\mathbf{A_1}$};
    \draw[thick] (A1.west) -- (2.0, 4.65);

    % Etiqueta A2
    \node[] (A2) at (4.5, 4.2) {$\mathbf{A_2}$};
    \draw[thick] (A2.east) -- (6.2, 4.2);

\end{tikzpicture}
